\textbf{1.} Deux raisons. D'une part, une mesure au rapporteur est
approchée (à un ou deux degrés près) : elle ne distingue pas $540^\circ$ de
$540{,}5^\circ$. D'autre part, on ne trace qu'un nombre fini de pentagones,
alors que l'affirmation porte sur \emph{tous} les pentagones convexes.

\textbf{2.} Soit $ABCDE$ un pentagone convexe. Les diagonales $[AC]$ et $[AD]$
sont intérieures au pentagone (exercice 0.3) et le découpent en trois
triangles $ABC$, $ACD$ et $ADE$. Les angles de ces trois triangles
reconstituent exactement les angles du pentagone : l'angle en $A$ est
partagé en trois, ceux en $C$ et en $D$ en deux, et ceux en $B$ et $E$ sont
des angles de triangles. La somme des angles d'un triangle valant
$180^\circ$, la somme cherchée vaut $3 \times 180^\circ = 540^\circ$.

\emph{Remarque.} Le même découpage, à partir d'un sommet, donne
$(n - 2) \times 180^\circ$ pour un polygone convexe à $n$ côtés.
\hfill $\blacksquare$
